\section{Resource bounds, lexicographic descent, heights, and complete lists}
\label{sec:resources}

Constructive formalization often fails not because the local step is difficult, but because an informal
proof says ``the process terminates'' or ``there are only finitely many possibilities'' without providing a
search cutoff.  iPRA needs a primitive-recursive resource.  This section gives two reusable conversions:
fixed-dimensional lexicographic descent to numerical fuel, and explicit height bounds to complete finite
lists.

\subsection{Mixed-radix rank for resetting counters}

Let a state contain a fixed tuple
\[
 a=(a_0,\ldots,a_{r-1}),
 \qquad 0\le a_i\le M_i.
\]
Coordinate \(0\) is most significant.  A strict lexicographic descent may decrease \(a_i\) while resetting
all lower-priority coordinates \(a_j\) with \(j>i\) to arbitrary values below their caps.  Define
\[
 \rank(a)=
 \sum_{i=0}^{r-1} a_i\prod_{j=i+1}^{r-1}(M_j+1).
\]
Equivalently, \(\rank\) is the mixed-radix packing operation.

\begin{lemma}[Lexicographic rank]
If \(b\) is strictly lexicographically smaller than \(a\), with both vectors inside the caps, then
\(\rank(b)<\rank(a)\).
\end{lemma}

\begin{proof}
Let \(i\) be the first coordinate at which they differ.  Then \(b_i<a_i\) and \(b_j=a_j\) for \(j<i\).
The decrease contributed at coordinate \(i\) is at least
\[
  W_i=\prod_{j=i+1}^{r-1}(M_j+1).
\]
The maximum possible increase in all lower coordinates is
\[
 \sum_{j=i+1}^{r-1}M_j
       \prod_{\ell=j+1}^{r-1}(M_\ell+1)
 =W_i-1,
\]
a standard mixed-radix identity proved by finite induction.  Hence the net rank change is at most \(-1\).
\end{proof}

\begin{corollary}[Fuel from fixed-dimensional descent]
A run that strictly decreases the capped lexicographic tuple has length at most
\[
  1+\rank(a^{0})
  \le \prod_{i<r}(M_i+1).
\]
\end{corollary}

For fixed \(r\), the rank, products, and comparison are primitive recursive, and the proof is formalizable
by quantifier-free induction on the finite sum/product identities.  The algorithm can therefore execute for
the displayed fuel and use the invariant ``every nonterminal step lowers rank.''

\subsection{Fixed versus input-dependent dimension}

The phrase ``lexicographic induction'' is ambiguous.  When the number of coordinates \(r\) is fixed by the
algorithm or generated module, mixed-radix rank gives an elementary primitive-recursive bound.  When \(r\)
is part of the input, the program must manipulate a list of counters and caps.  List packing, products, and
rank are still primitive recursive, but a proof may require bounds on code growth and on the length of the
list.  More sophisticated ordinal notations can exceed the strength or intended complexity of iPRA.

The framework therefore records dimension explicitly:
\begin{itemize}
  \item a fixed-dimensional module may compile the rank expression directly;
  \item a variable-dimensional theorem must supply a list-coded rank function and a proved numerical fuel;
  \item an informal appeal to an unspecified well order is not accepted as executable fuel.
\end{itemize}
This distinction resolves the ``resetting counters'' difficulty in several resource-oriented drafts without
claiming a general termination theorem beyond the supplied rank.

\subsection{Finite state-space fuel}

A second generic source of fuel is a finite state space.  If a deterministic transition remains inside
\([0,S)\), then within \(S+1\) states either a terminal state appears or some state repeats.  If the
application proves that a repeated nonterminal state is impossible---for example because a separate measure
strictly decreases---then termination occurs within \(S\) steps.  If repetition is allowed, cycle detection
itself is a finite theorem and may be the desired output.

A state-space argument is often better than deriving a sharp runtime estimate.  For a finite graph search,
Boolean table update, automaton, or bounded algebraic reduction, the number of valid packed states is an
explicit product of field caps.

\subsection{Height bounds and bounded enumeration}

Let \(P(x,z)\) be a decidable predicate on coded objects, and let \(H(x)\) be a height cutoff such that
\[
  \forall z\bigl(P(x,z)\to z\le H(x)\bigr).
\]
Define
\[
  L(x)=\langle z\le H(x):P(x,z)\rangle,
\]
the bounded filter of all codes up to \(H(x)\).  The filter is primitive recursive and its invariant proves
soundness, duplicate-freeness, and completeness.

\begin{theorem}[Height-to-list conversion]
\label{thm:height-list}
A primitive-recursive height cutoff with an object proof of completeness yields a primitive-recursive
complete finite list of all \(P\)-objects.
\end{theorem}
\begin{proof}
Enumerate \(z=0,\ldots,H(x)\), appending \(z\) precisely when the Booleanized predicate is true.  The loop
invariant says that the current list contains exactly the satisfying codes below the current counter.  At
\(H(x)+1\), soundness and bounded completeness are immediate.  The height theorem turns unbounded
completeness into bounded completeness: any satisfying \(z\) is below the cutoff and hence occurs.
\end{proof}

Conversely, suppose a finite list \(L(x)\) is explicitly complete:
\[
  \forall z\bigl(P(x,z)\to z\in L(x)\bigr).
\]
Then \(H(x)=\max L(x)\), with \(0\) for the empty list, is a height cutoff.  Thus, for decidable coded
objects, \emph{a complete finite list and a height bound are constructively interconvertible} when the
corresponding completeness proof is supplied.

\subsection{Why a cardinality bound is not a height bound}

A theorem of the form ``there are at most \(N\) objects'' does not provide a code cutoff.  The obstruction is
already visible for singletons.  For each parameter \(h\), define
\[
  P_h(z)\;:\!\!\Longleftrightarrow\;z=h.
\]
Every fiber has cardinality one, but its unique code can be arbitrarily large.  A uniform cutoff independent
of \(h\) cannot follow from cardinality one.  More generally, a finite set may occur at arbitrarily large
heights.

This distinction is crucial in Diophantine geometry.  A bound on the number of rational points, local
sections, or exceptional subspaces does not automatically yield an algorithm to enumerate them.  An actual
height bound, coefficient envelope, or complete list producer is stronger data.

\begin{warning}[No hidden finite choice]
A proof that for every input there exists some finite list covering all objects does not by itself define a
primitive-recursive function returning such a list.  The theorem must either supply a producer or remain a
boundary-conditional statement with the list as explicit input.
\end{warning}

\subsection{Finite envelopes}

A useful intermediate notion is a finite envelope.  Suppose every desired object is represented by an
auxiliary code \(a\) and a witness \(z\), and there is a bound \(A(x)\) on auxiliary codes plus a bound
\(H(x,a)\) on witnesses for each auxiliary code.  Then enumerate all pairs
\[
 a\le A(x),\qquad z\le H(x,a)
\]
and filter by the verifier.  This converts pointwise finite searches into one uniform complete list.  It is
the constructive content behind coefficient envelopes, finite covering families, bounded resultant
systems, and finite lists of local branches.

The envelope itself may be the missing deep theorem.  The framework makes its exact strength visible: a
finite cardinality assertion for each auxiliary fiber is insufficient unless it comes with a height or
list bound.

\subsection{Resource contracts}

A domain node should expose resource data through a contract such as
\[
  \operatorname{ValidInput}(x)\wedge\operatorname{ResourceCert}(x,r)
  \to
  \operatorname{RunCorrect}(x,r,F(x,r)).
\]
The resource certificate may contain caps, a rank function, coefficient heights, denominators, a gap, or a
finite state-space size.  A separate producer theorem may construct \(r\).  Keeping these layers separate
allows downstream proofs to be complete even when the best known global resource producer is conditional.

\subsection{Complexity accounting}

Primitive recursiveness alone can hide enormous growth.  The framework therefore records at least symbolic
resource expressions:
\begin{itemize}
  \item search bound and list length;
  \item loop fuel;
  \item demand cap and closure iterations;
  \item number of graph nodes and dependency uses;
  \item code-height growth under constructors;
  \item rational numerator and denominator bounds.
\end{itemize}
These bounds are not needed for logical soundness, but they prevent a proof from calling an operation
``finite'' while leaving its finite range unspecified.  They also guide later optimizations and make it
possible to compare alternative embeddings.
